\section{Conclusions and Future Work}\label{sec:conclusions}

This paper has presented \textbf{MemoizeLib}, an annotation-based memoization library whose design is shaped around three explicit goals: \emph{usability}, through an abstraction boundary that collapses the full correctness surface of a memoization refactoring behind a single annotation; \emph{automation-friendliness}, through an API deliberately designed to serve as the drop-in refactoring target of automated detectors and refactoring tools that have historically stopped at the detection step; and \emph{cross-platform reach}, through a single zero-dependency runtime module that serves both Android and the plain \jvm from identical semantics exposed through two parallel front-ends.

The library's runtime implements \numPolicies eviction policies, a dual-map \ttl subsystem whose disabled-state hot path is allocation-free, an auto-monitor that disables caches whose observed hit rate makes them counter-productive, and an embedded observability layer that incurs only a single volatile read per call site at its default-off level. Two sibling annotations---\code{@CacheInvalidate} and \code{@InvalidateCacheEntry}---express both bulk and single-row invalidation contracts declaratively, so that automated refactorers can encode the full invalidation surface of a mutating method without emitting any imperative code.

Alongside the library we have released an open Jetpack Benchmark harness that exercises \numAlgorithms compute-intensive, overlapping-subproblem algorithms in a baseline implementation and three memoized variants, one per eviction policy, across \numBenchmarkSuites benchmark suites that sweep per-algorithm speedup, per-policy hot-path cost, and hit-rate break-even. The library, the harness, the Sphinx documentation, and the analysis template build and link end-to-end as reproducible Gradle projects; the device measurement campaign is part of the replication package.

From the analysis grounded in the asymptotic complexity of the benchmarked algorithms and the implementation details of each policy, three conclusions follow. First, memoization converts every exponential naive recurrence in the benchmark set into a polynomial cached one, producing an algorithmic speedup that is expected to dominate the per-call instrumentation cost for any non-trivial input size. Second, the three eviction policies order as \fifo~$<$~\lru~$<$~\lfu in hot-path cost under no-eviction workloads, with the ordering reversing partially under eviction pressure on skewed distributions, where \lfu's frequency signal outperforms \lru's recency signal. Third, the hit-rate break-even point for recursion-heavy workloads lies well below the default \code{autoMonitor} threshold, which provides a conservative runtime safety net for automated refactoring tools whose static analysis may mis-estimate hit rate.

\paragraph{Future work.}
Several directions follow naturally from this work.

\begin{enumerate}[leftmargin=*,nosep]
    \item \textbf{A W-TinyLFU policy.} Caffeine's admission policy based on a count-min sketch is known to out-hit both \lru and \lfu on a wide range of workloads~\cite{einziger2017tinylfu}. The \code{MemoCache} interface is ready to accept it as a fifth implementation, at the cost of a small per-dispatcher sketch allocation.
    \item \textbf{Class-scoped caches.} The \code{CacheScope.CLASS} option currently declared on the annotation is not yet implemented in the transform. Completing it would enable a single cache to be shared across all instances of a pure-function-bearing class, eliminating redundant per-instance caches.
    \item \textbf{\code{onTrimMemory} integration.} The \code{MemoCacheManager} could register for Android's \code{ComponentCallbacks2.onTrimMemory} callback and proactively shrink caches under memory pressure, rather than waiting for GC to reclaim unreferenced dispatcher entries.
    \item \textbf{Detector integration.} The natural follow-up study is to couple MemoizeLib with an automated detector such as iMMutator~\cite{immutator} or an \llm-driven refactoring tool, and to evaluate the end-to-end pipeline---from detected candidate to merged \pr---in terms of both precision (fraction of suggestions accepted by maintainers) and delivered speedup (measured through the same Jetpack Benchmark harness on the modified code).
    \item \textbf{A broader device campaign.} Extending the measurement campaign beyond the Pixel family to additional vendors and to older SoCs would test whether the policy orderings hold across a wider slice of the Android fleet, following the cross-device validation pattern established in~\cite{empirical_rua}.
\end{enumerate}

The library, the Sphinx-generated documentation, the benchmark harness, and the \LaTeX{} source for this paper are released together as a reproducible research artefact, intended to serve as the downstream target for future detection tools, refactoring studies, and policy implementations.
